\documentclass{article}
\usepackage{amsmath,amssymb,xcolor,dashrule,iftex,enumitem}
\begin{document}
\section{Text formatting}
The Home ribbon writes ordinary TeX commands: \textrm{Roman},
\textsf{sans serif}, \texttt{monospace}, \textbf{bold}, and \textit{italic}.
\textcolor[HTML]{285E9C}{Text colors} remain part of the document source.

{\small Small text}, {\normalsize normal text}, and {\Large large text}.

\begin{flushleft}Left aligned paragraph, including inline math $x^2+1$.\end{flushleft}

\begin{center}Centered paragraph with \textbf{bold words}.\end{center}

\begin{flushright}Right aligned paragraph.\end{flushright}

Normal paragraphs use the TeX document's default alignment.

\begin{enumerate}[label=(\alph*),start=3,nosep,leftmargin=2em]
\item A compact list starting at (c).
\item Change the label, starting number or spacing in its source options.
\end{enumerate}

\noindent Left\hfill Centre\hfill Right

\noindent Contents\dotfill Page 1

\noindent Signature\hrulefill Date\hrulefill

The selected engine is
\ifXeTeX Xe\protect\LaTeX\else
  \ifLuaTeX Lua\protect\LaTeX\else pdf\protect\LaTeX\fi
\fi{}; inactive branches remain in the source.

\paragraph{Run-in heading}The body continues on the same line in the normal font.

\subparagraph{Indented heading}Only the heading's first line is indented.
Longer text wraps at the normal paragraph margin.

\setcounter{secnumdepth}{5}
\paragraph{Numbered heading}\label{par:numbered}Numbering depth controls which heading levels receive numbers.
\subparagraph{Numbered detail}This is a child of paragraph~\ref{par:numbered}.
\setcounter{secnumdepth}{3}

\section{Combining styles}
\textsf{\textit{\textbf{Bold italic sans serif}}},
\textrm{\textsc{Small capitals}}, and
\textsl{\textsc{Slanted small capitals}} combine family, weight and shape.
Select only part of a styled phrase. The ribbon should show the active styles;
clicking an active style removes it and preserves the selection. Change the
family list from Roman to Sans serif, then inspect the source.

\section{Sizes, color and literal characters}
{\footnotesize Footnote size}, {\small small}, {\normalsize normal},
{\large large}, and {\LARGE larger text} use named sizes.
{\fontsize{13pt}{16pt}\selectfont This sentence uses a custom size and baseline.}
Use Home, Font, Custom to set a number and unit, then undo.

\textcolor[HTML]{A33D30}{A colored phrase with \textbf{bold words} inside it.}
Select individual characters across the formatting boundary. Restore default
color with the triangle beside Text color, or use Clear formatting.

Literal characters: \% \$ \& \# \_ \{ \} \textbackslash{}.
Text angle brackets: \textless{}tag\textgreater{}; math relations: $a<b>c$.
The logos are \TeX{} and \LaTeX{}. This line ends explicitly.\newline
The next line still belongs to the same paragraph.

\section{Boxes and local declarations}
Two top-aligned minipages keep their own line width:

\noindent\begin{minipage}[t]{0.46\linewidth}
Text and $x^2+1$ in a narrow column.
\par\hrule height 0.6pt
A solid rule separates these lines.
\end{minipage}\hspace{0.04\linewidth}%
\begin{minipage}[t]{0.46\linewidth}
A second column with a dashed rule.
\par\noindent\hdashrule{\linewidth}{0.6pt}{3pt 2pt}
\end{minipage}

Text and mathematics can share a frame: \fbox{Text $x^2+\frac{1}{y}$}.
Compare \framebox[8em][c]{Centered $a+b$} and \mbox{Unbroken text $n^2$}.
Click inside a box to open the Box tab; Box settings changes its contents,
width or alignment. Edit box source expands its code in place.

Fixed-height boxes support outer and inner alignment:
\parbox[t][3em][t]{.35\linewidth}{Top-aligned text\\on two lines.}\quad
\parbox[t][3em][b]{.35\linewidth}{Bottom-aligned text\\on two lines.}

{\small\bfseries Local size and weight declarations apply only inside this group.}
Hover over a group or command to reveal its dashed outline. The plus button
opens editable source; the minus collapses it. These controls belong to the editor overlay.

% Select this comment and use Ctrl+/ to turn it into ordinary text.
\section{Selection and scrolling practice}
Drag across several lines and paragraphs, then scroll using the wheel while
leaving the caret in the document. The selection and viewport should remain
where you put them. Double-click a word, extend with Shift+arrows, copy it, and
compare the selection in the full source editor using Ctrl+Shift+U.

This second paragraph provides another selection endpoint. Ctrl+L shows a
gray source box for a selection; toggling it back keeps the source unchanged.

\section{Literal text}
In PasTeX, click inside either example and type. Enter inserts a literal line
break; braces, percent signs and backslashes keep their original characters.
Undo and redo work as in the body. The source remains ordinary \LaTeX.

\begin{verbatim}
\frac{x}{y}  $literal delimiters$  % literal percent
    indented text with {braces} and a backslash: \
two  spaces    four spaces
\end{verbatim}

The starred environment makes spaces visible in both previews:
\begin{verbatim*}
one space  two spaces   three spaces
\unknown{literal command}
\end{verbatim*}

Use inline \verb|\verb| for a short literal fragment, for example
\verb|C:\work\notes.tex|. Unsupported literal environments from other packages
keep their complete editable source box.
\section{Footnotes and margin notes}
This sentence has an editable footnote\footnote{A note can contain
\textbf{formatted text} and mathematics such as $x^2+1$.}.
The next mark and its text are separate commands\footnotemark.
\footnotetext{This note belongs to the preceding mark.}
A short margin note\marginpar{Review this paragraph.} demonstrates placement
outside the text column. The compiled paged view uses the document's real layout.
\end{document}
